\section{Translation invariance and geometric representation}

\begin{lemma}[Translation invariance]\label{lem:translation}
For any $\boldsymbol c\in\R^d$ and $\boldsymbol\alpha$ satisfying $\boldsymbol t^\T\boldsymbol\alpha=0$,
\[
\sum_{j,k}\alpha_j\alpha_k t_jt_k\boldsymbol x_j^\T\boldsymbol x_k
=
\sum_{j,k}\alpha_j\alpha_k t_jt_k(\boldsymbol x_j-\boldsymbol c)^\T(\boldsymbol x_k-\boldsymbol c).
\]
\end{lemma}
\begin{proof}
Bilinearity gives
\[
\sum_{j,k}\alpha_j\alpha_kt_jt_k\boldsymbol x_j^\T\boldsymbol x_k
=\left\lVert\sum_j\alpha_jt_j\boldsymbol x_j\right\rVert^2.
\]
Moreover,
\[
\sum_j\alpha_jt_j\boldsymbol x_j
=\sum_j\alpha_jt_j(\boldsymbol x_j-\boldsymbol c)
+\left(\sum_j\alpha_jt_j\right)\boldsymbol c
=\sum_j\alpha_jt_j(\boldsymbol x_j-\boldsymbol c).
\]
\end{proof}

Take $\boldsymbol c=\boldsymbol\mu_*$ and define the signed centered Gram matrix $\boldsymbol G=(G_{jk})$ by
\[
G_{jk}=t_jt_k(\boldsymbol x_j-\boldsymbol\mu_*)^\T(\boldsymbol x_k-\boldsymbol\mu_*).
\]
Then
\begin{equation}\label{eq:centered-dual}
L(\boldsymbol\alpha)=\one^\T\boldsymbol\alpha-\frac12\boldsymbol\alpha^\T\boldsymbol G\boldsymbol\alpha,
\qquad \boldsymbol\alpha\in\cC.
\end{equation}
Its deterministic leading term is
\begin{equation}\label{eq:Gtilde}
\widetilde G_{jk}=
\begin{cases}
\Delta/4+\delta_{jk}\tr(\boldsymbol\Sigma_i),&j,k\in I_1\cup I_2,\ j\in I_i,\\
D_k\delta_{jk},&j,k\in O,\\
0,&\text{otherwise},
\end{cases}
\end{equation}
where $\delta_{jk}$ is the Kronecker delta.

\begin{lemma}[Geometric representation with outliers]\label{lem:geometry}
Assume (A1)--(A2), and let $\boldsymbol R=\boldsymbol G-\widetilde{\boldsymbol G}$. Then
\[
\epsilon_1:=\max_{(j,k)\notin O\times O}{|R_{jk}|\over\Delta_*}=\op(1),
\qquad
\epsilon_2:=\max_{k\in O}{|R_{kk}|\over D_k}=\op(1).
\]
\end{lemma}
\begin{proof}
Put $\boldsymbol u=-\boldsymbol\mu/2$. For every inlier,
$t_j(\boldsymbol x_j-\boldsymbol\mu_*)=\boldsymbol u+t_j\boldsymbol\varepsilon_j$,
whereas for $k\in O$,
$t_k(\boldsymbol y_k-\boldsymbol\mu_*)=\boldsymbol\nu_k+\boldsymbol\eta_k$.
Also $\lVert\boldsymbol u\rVert^2=\Delta/4$.

For distinct inliers $j\in I_i$ and $k\in I_l$,
\[
R_{jk}=t_j\boldsymbol u^\T\boldsymbol\varepsilon_j
+t_k\boldsymbol u^\T\boldsymbol\varepsilon_k
+t_jt_k\boldsymbol\varepsilon_j^\T\boldsymbol\varepsilon_k.
\]
The three terms have mean zero. Their second moments are controlled by
$\boldsymbol\mu^\T\boldsymbol\Sigma_i\boldsymbol\mu/4$,
$\boldsymbol\mu^\T\boldsymbol\Sigma_l\boldsymbol\mu/4$, and
$\tr(\boldsymbol\Sigma_i\boldsymbol\Sigma_l)$; Chebyshev's inequality and (A1) therefore give $R_{jk}/\Delta_*=\op(1)$.
For an inlier diagonal,
\[
R_{jj}=2t_j\boldsymbol u^\T\boldsymbol\varepsilon_j+
\{\lVert\boldsymbol\varepsilon_j\rVert^2-\tr(\boldsymbol\Sigma_i)\},
\]
and (A1) gives the same conclusion.

For $j\in I_i$ and $k\in O$,
\[
R_{jk}=G_{jk}=-{\boldsymbol\mu^\T\boldsymbol\nu_k\over2}
-{\boldsymbol\mu^\T\boldsymbol\eta_k\over2}
+t_j\boldsymbol\varepsilon_j^\T\boldsymbol\nu_k
+t_j\boldsymbol\varepsilon_j^\T\boldsymbol\eta_k.
\]
Assumption (A2) makes the deterministic term negligible relative to $\Delta_*$ and bounds the second moments of the other terms. For $k\ne l$ in $O$,
$R_{kl}=\boldsymbol\nu_k^\T\boldsymbol\nu_l+\boldsymbol\nu_k^\T\boldsymbol\eta_l+\boldsymbol\nu_l^\T\boldsymbol\eta_k+\boldsymbol\eta_k^\T\boldsymbol\eta_l$, which is again $\op(\Delta_*)$ by (A2). Finally,
\[
R_{kk}=2\boldsymbol\nu_k^\T\boldsymbol\eta_k+
\{\lVert\boldsymbol\eta_k\rVert^2-\tr(\boldsymbol\Sigma_{0k})\}=\op(D_k).
\]
Because $N$ is fixed, a union bound over the finitely many entries proves both maxima.
\end{proof}
